% Draft 2026-11-03: official setup, only the model differs. FORGE 2027, IEEEtran, double-anonymous: no repo URLs, names or cluster identifiers.
\section{Method}
\label{sec:method}

\subsection{Setup}
We run FrontierMemory's public release~\cite{frontiermemory} through its own commands (\texttt{generate-memory}, \texttt{run-tasks}, \texttt{score}, \texttt{bundle}) and task containers, on Docker.
The leaderboard's default generator and loader are not public, so we use Hermes, a published leaderboard row whose generator and task agent are both in the public repository at a pinned commit.
All 15 tasks are run on corpus v1.1, as used on the leaderboard, although the public code still pins v1.0.
A run is one fresh memory generation under the official 6-hour limit followed by 3 attempts per task; we do 3 runs (135 trials), plus a no-memory control of the same agent with 3 attempts per task (45 trials).
Official timeouts apply.

\subsection{What differs from official runs}
\paragraph{Model} Generation and tasks both use Qwen3.8-27B (4-bit) served by Ollama, at its full 262{,}144-token context. A thin adapter, loaded through the official custom-harness hook, changes only the model URL. Effort is requested at the highest level; Hermes sends a model it doesn't recognize \emph{xhigh}.
\paragraph{Docker host} The cluster has no Docker, so the official Docker Engine runs in a virtual machine, with the model on the host GPUs. The official network probe and reference solutions behave as the maintainers report (\S\ref{sec:method-validity}).
\paragraph{One task repair} Cookie-store's verifier cannot be built as published since a dependency released on Nov 1 (\S\ref{sec:results}). We add the version its own test script requires, one line in the verifier's build.
\paragraph{Judge} Three APEX tasks are graded by an LLM judge, documented as GPT-5.6 Sol. \todo{We use Claude Opus 5.5, or regraded results if the documented judge becomes available.} It reproduced the reference scores exactly; agreement on wrong answers is untested.
\paragraph{Image} Word2Vec's image installs from a floating branch, so we pre-build it once, unmodified, and reuse it for every trial.

\subsection{Validity}
\label{sec:method-validity}
Before any counted trial, every task's reference solution was run and matched the maintainers' reported scores: 10 tasks at 1, one at about 0.83 and 4 at 0.
The official network probe shows only the model host reachable during generation.
One small task then ran end to end at full settings, not counted.

\subsection{Analysis}
Design and analysis were preregistered before any counted trial.
The primary outcome is strict pass.
We test memory against no memory by exact one-sided paired sign-flip over tasks, pooling a task's memory attempts across runs, and report effect sizes with exact intervals.
With this test, $p<.05$ needs at least 5 tasks moving the same way with none opposing, so a null is not evidence of no effect.

\subsection{Audit}
For each link we report what we can establish and how: from source and documentation, by measurement, or by inference.
\begin{description}
  \item[Corpus] Per-task coverage, from prompt similarity and from repository keywords, for both corpus versions.
  \item[Generation] What each memory contains and how much of the corpus the generator read, from its trace.
  \item[Harness] The rule each harness uses to load memory, and the share of a memory it would show.
  \item[Agent] Whether trials open memory files, from traces.
  \item[Grader] Reference-solution behaviour, judge identity, handling of judge failure, and build reproducibility.
\end{description}
Findings about construction apply to any model; those about gains from memory, and generation quality, apply to this model only.
